\documentclass[11pt]{amsart}
\usepackage[T1]{fontenc}
\IfFileExists{lmodern.sty}{\usepackage{lmodern}}{}
\usepackage{amssymb,amsmath}
\usepackage[hidelinks]{hyperref}
\newtheorem{theorem}{Theorem}
\newtheorem{lemma}[theorem]{Lemma}

% Generated by erdos-engine 0.1.0 on 2026-10-06.
% TODO(human): every passage marked TODO must be checked and rewritten by the
% author before this goes anywhere.

\title{A certified computational verification for Erd\H{o}s Problem \#458}
\author{AUTHOR NAME (fill in)}
\address{Affiliation (fill in)}
\email{email@example.org}
\date{2026-10-06}
\subjclass[2020]{11A41 (secondary 11Y11)}

\begin{document}
\begin{abstract}
Erd\H{o}s and Graham asked whether $\operatorname{lcm}(1,\ldots,p_{k+1}-1)<p_k\operatorname{lcm}(1,\ldots,p_k)$ for every $k$, where $p_k$ is the $k$-th prime. We verify this for all $k$ with $p_{k+1}\le 10^{21}$. The proof is self-contained: it locates every prime power with exponent at least $3$ below the bound exactly and exhibits a Pocklington-certified prime between every pair of consecutive prime squares, so it does not rely on tables of maximal prime gaps. Every claim is backed by a certificate that was re-validated by an independent implementation.
\end{abstract}
\maketitle

\section{The problem}
Erd\H{o}s Problem \#458 \cite{erdosproblems458} asks the following.
\begin{quote}
Let $[1,\ldots,n]$ denote the least common multiple of $\{1,\ldots,n\}$ and $p_k$ the $k$-th prime. Is it true that, for all $k\geq 1$, \[[1,\ldots,p_{k+1}-1]< p_k[1,\ldots,p_k]?\]
\end{quote}
At the time of writing it is listed as \emph{falsifiable} in the community
database \cite{ErdosDB}: it is open, but a counterexample, if one exists, is a finite object.

\subsection*{Prior work}
% TODO(human): the sentence below was assembled automatically from the
% problem's forum thread.  Read the thread and the literature, and rewrite.
The forum thread for the problem contains 3 comments. The largest verification bound reported there is $10^{20}$ (bhowerton, erdosproblems.com forum, 12 Jun 2026; that verification relies on prime-gap records above 4e18). The present computation goes beyond that bound.

\section{Method}

\paragraph{Reduction.} For consecutive primes $p<r$ the primes dividing
$[1,\ldots,r-1]/[1,\ldots,p]$ are exactly the bases $q$ of prime powers
$q^a$ ($a\ge 2$) in the open interval $(p,r)$, each contributing one factor
$q$ per such power. Hence the inequality for the gap $(p_k,p_{k+1})$ is
equivalent to $\Pi(p_k,p_{k+1})<p_k$, where $\Pi(p,r)=\prod_{q^a\in(p,r)}q$.
(This reformulation also appears in the forum discussion of the problem.)

\paragraph{Certificate.} We prove the statement for all $k$ with
$p_{k+1}\le X$ without using any table of maximal prime gaps:
\begin{enumerate}
\item[(A)] For every prime power $h=q^a\le X$ with $a\ge 3$ we determine the
enclosing gap $(p,r)\ni h$ exactly (deterministic primality testing of
every integer between), list \emph{all} prime powers inside it, squares
included, and verify $\Pi(p,r)<p$.
\item[(B)] For every pair of consecutive primes $q<q'$ with $q\le\sqrt{X}$ we
exhibit a prime in $(q^2,q'^2)$. Consequently no prime gap contains two
prime squares.
\item[(C)] A gap with $r\le X$ containing no prime power of exponent $\ge3$
thus contains at most one prime square $q^2$, so $\Pi(p,r)\le q<p$ (there
is a prime in $(q,q^2)$ by Bertrand's postulate).
\end{enumerate}

\paragraph{Primality.} All primality decisions use either the Miller--Rabin
test with the first 13 prime bases, which is deterministic below
$3.317\cdot10^{24}$ \cite{SorensonWebster2017}, or a Pocklington--Lehmer
certificate \cite{Pocklington1914,BLS75}: the witnesses in (B) are chosen of
the form $N=Pm+1$ with $P$ one of the 32 largest primes $\le q$ and $m$ even,
so that $F=2P>\sqrt N$ is a completely factored divisor of $N-1$ and a single
base $a$ with $a^{N-1}\equiv1$, $\gcd(a^{(N-1)/P}-1,N)=\gcd(a^{(N-1)/2}-1,N)=1$
proves $N$ prime. Compositeness is only ever asserted from an explicit failed
strong-probable-prime test, which is a proof.

\paragraph{Independent checking.} The search is implemented in C (two-limb
Montgomery arithmetic). A separate pure-Python implementation, sharing no
code with it, (i) re-verifies the literal lcm inequality and the reduction on
an initial range, (ii) recomputes part (A) in full and compares a SHA-256
digest of all gap records, and (iii) recomputes randomly sampled chunks of
part (B), comparing a 64-bit FNV-1a digest of the entire witness sequence.
The witness rule is a pure function of $(q,q')$, so agreement of digests means
both implementations produced identical witnesses.


\section{Results}

\begin{theorem}
For every $k$ with $p_{k+1}\le 10^{21}$,
\[[1,\ldots,p_{k+1}-1]<p_k\,[1,\ldots,p_k].\]
\end{theorem}

Part (A) located $683532$ prime powers with exponent $\ge3$, lying in
$683531$ distinct prime gaps, of which $5$ contain more than one prime
power. Part (B) exhibited $1367199811$ witnesses (one per prime $q\le 31622776601$),
in $3163$ chunks; $3782301$ of them needed the Miller--Rabin fallback
rather than a Pocklington certificate. The square scan used about
$6.2$ CPU-hours. The literal lcm inequality was additionally checked
directly for $p_{k+1}\le 3000$.
Limitation: the independent re-implementation recomputed part (A) in full
but part (B) only for a random sample of chunks together with the first
and last (Section~4); the coverage of (B) is confirmed by comparing the
number of witnesses with $\pi(\lfloor\sqrt X\rfloor)$ computed by a
sieve-free method, while correctness of the witnesses outside the sample
rests on the C implementation and its per-witness primality proofs.

The gaps with the largest ratio $\Pi(p,r)/p$ are all small:
\begin{center}
\begin{tabular}{lll}
gap $(p,r)$ & prime powers inside & $\Pi(p,r)$ \\ \hline
$(7,\,11)$ & $2^{3} \cdot 3^{2}$ & $6$ \\
$(23,\,29)$ & $3^{3} \cdot 5^{2}$ & $15$ \\
$(113,\,127)$ & $5^{3} \cdot 11^{2}$ & $55$ \\
$(13,\,17)$ & $2^{4}$ & $2$ \\
$(31,\,37)$ & $2^{5}$ & $2$ \\
$(79,\,83)$ & $3^{4}$ & $3$ \\
\end{tabular}
\end{center}


\section{Verification and reproducibility}
The certificate produced by the search was re-validated by an independent
implementation (``p458-independent-python'') that shares no code with the search:
\begin{itemize}
\item literal lcm inequality re-verified for p\_(k+1) <= 3000
\item reduction (product criterion) agrees with literal lcm on that range
\item (A) 683532 prime powers / 683531 gaps recomputed independently; digest matches
\item (B) 3163 chunks tile [2, 31622776601] with 1367199811 primes, no failures reported
\item (B) witness count equals pi(31622776601) = 1367199811, computed by the Lucy\_Hedgehog recursion (no sieve): every prime q <= sqrt(X) has a witness
\item 26 chunks (including the first and the last, the others chosen at random) recomputed from scratch and reproduced exactly: 11,182,376 witnesses in total
\end{itemize}
The raw certificate (\texttt{result.json}, SHA-256
\texttt{08344b04871cb70f cc2b50e4a69dc9b2 a7f39ac6a4eb6449 b7adf266b88258cc}) and the check report are deposited
with the software. % TODO(human): add repository URL and Zenodo DOI.

\section*{Use of AI tools}
The search and checking software was written with the assistance of an AI coding assistant (Claude, Anthropic). All numerical claims come from executed runs whose certificates were re-validated by an independent implementation; the author(s) reviewed the code and the mathematics and take full responsibility for the content.

\nocite{ErGr80,SorensonWebster2017,OeSHP2014,Pocklington1914,BLS75,erdosproblems458}
\bibliographystyle{amsplain}
\bibliography{refs}
\end{document}
